% Related Work — NeurIPS framing.
% Five positioning blocks:
%   (a) document KIE landscape with the structural-verification gap;
%   (b) verifier-augmented generation in NLP / reasoning;
%   (c) neuro-symbolic learning with hard constraints;
%   (d) arithmetic-consistency post-processing in invoice extraction —
%       the closest published precedent for FOCUS-Σ I3;
%   (e) receipt and document-AI benchmarks.
% Citation count: 38 distinct works (NeurIPS norm: 30–60).

\section{Related Work}
\label{sec:related}

\paragraph{Document KIE without verification.}
End-to-end document understanding has converged on two architectural
families.  Generative models — DONUT~\cite{kim2022donut},
Dessurt~\cite{davis2022dessurt}, Pix2Struct~\cite{lee2023pix2struct},
UDOP~\cite{tang2023unifying}, DocLLM~\cite{wang2024docllm}, and
LayoutLLM~\cite{luo2024layoutllm} — autoregressively decode the
structured output from the page image, optionally conditioned on OCR
tokens.  Layout-aware encoders fuse text, image, and 2-D position via
masked language / image objectives:
LayoutLM~\cite{xu2020layoutlm,xu2021layoutlmv2,huang2022layoutlmv3},
LayoutXLM~\cite{xu2021layoutxlm}, BROS~\cite{hong2022bros},
TILT~\cite{powalski2021tilt}, DocFormer~\cite{appalaraju2021docformer},
ERNIE-Layout~\cite{peng2022ernielayout}, and
StrucTexT~\cite{li2021structext}.  Pipeline approaches retain a
modular detect/read/assign decomposition with a text-line detector
(YOLOv8~\cite{jocher2023yolov8}, DBNet~\cite{liao2020dbnet},
PaddleOCR~\cite{paddleocr}, Tesseract~\cite{smith2007tesseract}),
a text recogniser (TrOCR~\cite{li2023trocr}), and either a learned
field-assignment head — exemplified by PICK~\cite{yu2020pick} and our
FOCUS-T baseline — or a rule-based assigner.  Across both families,
the cited systems emit one prediction per field per receipt without
any consistency check against the receipt's own structural
evidence.  Our contribution is a verification layer that wraps any
of these architectures.

\paragraph{Verifier-augmented generation in NLP.}
Pairing a generator with a downstream verifier has become a standard
device for chain-of-thought tasks.  Cobbe et al.~\cite{cobbe2021training}
train a verifier on GSM8K to score solution candidates;
Lightman et al.~\cite{lightman2023verify} contrast process- and
outcome-supervised reward models on PRM800K;
Madaan et al.~\cite{madaan2023selfrefine} use the generator itself as
iterative critic; Chen et al.~\cite{chen2022codet} verify generated
code against generated tests; Wang et al.~\cite{wang2023selfconsistency}
marginalise across sampled chains to amplify a hidden consensus.
Tool-using models — Toolformer~\cite{schick2023toolformer},
PAL~\cite{gao2023pal} — outsource verification to an external
interpreter, which is closest in spirit to FOCUS-$\Sigma$ but
typically defers an entire arithmetic step rather than checking a
single emitted candidate against a structural identity.
FOCUS-$\Sigma$ is in this lineage but differs in two respects.
\textbf{(i)} The verification target is a single numerical value
bounded by a known arithmetic identity rather than a free-form claim
or program; we exploit this to give exact soundness and completeness
guarantees (Sec.~\ref{sec:focus_sigma}).
\textbf{(ii)} The verifier is itself rule-based (a bounded subset-sum
DP) rather than a learned head, which keeps it
architecture-agnostic and trivially deployable across upstream
extractors with no new gradient signal.

\paragraph{Neuro-symbolic learning with hard constraints.}
A complementary literature integrates symbolic constraints \emph{into}
the learner rather than wrapping its outputs.  DeepProbLog
\cite{manhaeve2018deepproblog} differentiates through probabilistic
logic programs; the Semantic Loss~\cite{xu2018semanticloss} penalises
violations of propositional constraints during training;
Reluplex~\cite{katz2017reluplex} statically verifies neural networks
against linear specifications.  These methods change the network's
loss surface or formal contract.  FOCUS-$\Sigma$ takes the
operationally cheaper path: the upstream extractor remains untouched,
and a domain axiom (the receipt-construction rule) is enforced as a
post-hoc test.  This trades the in-training generalisation benefits
of constraint learning for a deployment property — the verifier
ships to production without any joint optimisation — that matters in
the document-AI setting where pre-trained backbones are repeatedly
swapped.

\paragraph{Arithmetic post-processing in invoice extraction.}
The closest published precedent is in the invoice-extraction
literature, where heuristic checks of \texttt{subtotal + tax = total}
appear as a post-processing step.
Kim et al.~\cite{kim2022donut} cite ``arithmetic-consistency
post-processing'' as an ablation;
Hong et al.~\cite{hong2022bros} mention sub-total + tax sanity checks
in their post-process; commercial systems
(Microsoft Form Recogniser~\cite{azure-form-recogniser},
AWS Textract~\cite{aws-textract}) advertise total-field validation
without publishing their rule set.  These approaches are
\emph{keyword-anchored}: they fire only when the \textsf{SUBTOTAL}
and \textsf{TAX} keywords are both present and OCR'd correctly, which
on the canonical SROIE Task-3 split holds for only 38.3\% of
receipts.  FOCUS-$\Sigma$ adds the structural Identity $\mathrm{I}_3$
which fires on the full 87.6\% of receipts that satisfy the
itemisation hypothesis, regardless of which summary keywords the
printer chose to emit.

\paragraph{Subset-sum and combinatorial-verification methods.}
The bounded subset-sum we use is classical knapsack DP
(Karp~\cite{karp1972reducibility}; Bellman~\cite{bellman1957dynamic}).
Its application as a verifier in document AI is, to our knowledge,
new.  Closer methodologically is the program-by-construction
literature
(Solar-Lezama~\cite{solarlezama2008program};
Polozov \& Gulwani~\cite{polozov2015flashmeta}) that synthesises
domain-specific constraints from examples, but those systems
synthesise rules; FOCUS-$\Sigma$ uses a fixed rule (a domain axiom)
and verifies against it.

\paragraph{Faithfulness and calibration of structured predictors.}
The interpretability claim we make about the FOCUS-T attention map is
consistent with a literature that quantifies attribution faithfulness
through deletion / insertion AUCs~\cite{petsiuk2018rise} and
benchmarks rationale extraction~\cite{deyoung2020eraser}.
Calibration of the verifier-augmented predictor follows the standard
ECE / MCE / Brier toolkit~\cite{guo2017calibration,naeini2015ece,brier1950score};
Sec.~\ref{sec:limitations} discusses why these measurements are
preregistered rather than reported on the headline run.

\paragraph{Receipt and document benchmarks.}
SROIE~\cite{huang2019icdar} (626 receipts, 4 fields) is the standard
single-vendor-type benchmark.  CORD~\cite{park2019cord} (Korean
restaurant receipts, 30-field schema) is the second most-studied
receipt corpus.  WildReceipt~\cite{sun2021wildreceipt} extends to
in-the-wild imaging conditions but is rarely used as a fine-tune
target.  Form / page benchmarks adjacent to receipts include
FUNSD~\cite{jaume2019funsd}, RVL-CDIP~\cite{harley2015rvlcdip},
DocVQA~\cite{mathew2021docvqa}, and Kleister~\cite{stanislawek2021kleister};
they share the layout-driven extraction motif but lack the closed
arithmetic structure that makes FOCUS-$\Sigma$ applicable.  Our
unified text-to-text framing for receipts borrows from
T5~\cite{raffel2020t5} and BART~\cite{lewis2020bart}; the visual
backbone is Swin~\cite{liu2021swin}.  We evaluate on SROIE Task-3 and
CORD-v2; WildReceipt generalisation is flagged as future work
(Sec.~\ref{sec:limitations}).
